\begin{helper}[A proper initial segment has smaller Tashkinov-tree rank]
If \(S\) is a proper initial segment of a tree-sequence \(T\), then
\(S\) has smaller lexicographic rank \((\rho,\text{length})\) than
\(T\).
\end{helper}
\begin{proof}
Let \(q=\rho(T)\), where \(\rho\) is
\noderef{TashTreePathIndex}.  By
\noderef{TashTreeSuffixPath}, \(1\le q\le |E(T)|\), and every displayed
edge at or after position \(q\) joins its two consecutive displayed
vertices.  The same definition gives \(\rho(S)\le |E(S)|\).

If \(|E(S)|<q\), then
\(\rho(S)\le |E(S)|<q=\rho(T)\), so the first rank coordinate strictly
decreases.  Otherwise \(q\le |E(S)|\).  Since
\noderef{TreeSequenceInitialSegment} identifies every retained vertex
and exact edge of \(S\) with the corresponding entry of \(T\), all the
suffix-path equations of \(T\) from position \(q\) restrict to \(S\).
Consequently \(q\) is an admissible suffix start for \(S\), and the
minimality in \noderef{TashTreePathIndex} gives \(\rho(S)\le q\).  If
this inequality is strict, the first coordinate again decreases.  If
it is an equality, the first coordinates agree while the properness
hypothesis gives \(|E(S)|<|E(T)|\).  In every case the lexicographic
condition \noderef{TashTreeRankLess} holds.
\end{proof}
